% !TEX program = lualatex
% !TeX encoding = UTF-8
\PassOptionsToPackage{unicode}{hyperref}
\PassOptionsToPackage{bookmarksnumbered,bookmarksopen}{hyperref}
\documentclass[12pt,a4paper]{article}

% ============================================================
% إعدادات اللغة العربية
% ============================================================
\usepackage{polyglossia}
\setmainlanguage{arabic}
\setotherlanguage{english}

% خطوط عربية
\newfontfamily\arabicfont{Adobe Arabic}[Script=Arabic, Language=Arabic]
\newfontfamily\arabicfontsf{Adobe Arabic}[Script=Arabic, Language=Arabic]
\newfontfamily\arabicfonttt{Adobe Arabic}[Script=Arabic, Language=Arabic]
\newfontfamily\englishfont{Times New Roman}
\setmonofont{Latin Modern Mono}

% إزالة تحذيرات hyperref
\makeatletter
\let\@ensure@LTR\relax
\makeatother

% حزم الرياضيات والتنسيق
\usepackage{amsmath,amssymb,amsthm,mathtools,bm}
\usepackage{geometry}
\usepackage{enumitem}
\usepackage{siunitx}
\usepackage{booktabs}
\usepackage{array}
\usepackage{graphicx}
\usepackage{caption}
\usepackage{makecell}
\usepackage{pgfplots}
\usepackage{float}
\usepackage{tikz}
\usepackage{multicol}
\usepackage{microtype}
\usepackage{hyperref}
\usepackage{listings}
\usepackage{xcolor}
\usepackage{tcolorbox}
\usepackage{algorithm}
\usepackage{algorithmic}
\pgfplotsset{compat=1.18}
\usetikzlibrary{arrows.meta,positioning,shapes.geometric,fit,calc}

\geometry{margin=2.5cm}
\setlength{\emergencystretch}{3em}
\sloppy

% Theorem environments
\newtheorem{theorem}{نظرية}
\newtheorem{lemma}{ليما}
\newtheorem{corollary}{نتيجة طبيعية}
\newtheorem{definition}{تعريف}
\newtheorem{axiom}{بديهية}
\newtheorem{proposition}{اقتراح}
\newtheorem{remark}{ملاحظة}
\newtheorem{assumption}{افتراض}
\renewenvironment{proof}{\paragraph{برهان:}}{\hfill$\square$}

% Operators
\DeclareMathOperator{\Tr}{Tr}
\DeclareMathOperator{\Diff}{Diff}
\DeclareMathOperator{\Aut}{Aut}
\DeclareMathOperator{\Vol}{Vol}
\DeclareMathOperator{\Ric}{Ric}
\DeclareMathOperator{\Riem}{Riem}
\DeclareMathOperator{\Cov}{Cov}
\DeclareMathOperator{\supp}{supp}
\DeclareMathOperator{\diag}{diag}
\DeclareMathOperator{\sign}{sign}
\DeclareMathOperator{\dimension}{dim}
\DeclareMathOperator{\Div}{Div}
\DeclareMathOperator{\Grad}{Grad}
\DeclareMathOperator{\Hess}{Hess}
\DeclareMathOperator{\Ricci}{Ricci}
\DeclareMathOperator{\Scalar}{Scalar}

% Robust math shorthands
\DeclareRobustCommand{\alD}{\texorpdfstring{\ensuremath{\alpha_{\mathrm{D}}}}{alpha_D}}
\DeclareRobustCommand{\beI}{\texorpdfstring{\ensuremath{\beta_{\mathrm{I}}}}{beta_I}}
\DeclareRobustCommand{\gaR}{\texorpdfstring{\ensuremath{\gamma_{\mathrm{R}}}}{gamma_R}}
\DeclareRobustCommand{\UCS}{\texorpdfstring{\ensuremath{\mathcal{M}_{\mathrm{UCS}}}}{M_UCS}}

% YUCT-specific commands
\newcommand{\eff}{\mathrm{eff}}
\newcommand{\coord}{\mathrm{coord}}
\newcommand{\Yak}{\mathrm{Yak}}
\newcommand{\Dict}{\mathcal{D}}
\newcommand{\Mdict}{\mathcal{M}_{\Dict}}
\newcommand{\Ke}{K_{\eff}}
\newcommand{\Kemax}{K_{\eff,\max}}
\newcommand{\Keobs}{K_{\eff,\mathrm{obs}}}
\newcommand{\Keexp}{K_{\eff,\mathrm{exp}}}
\newcommand{\Kec}{K_{\eff,c}}
\newcommand{\Kcrit}{K_{\mathrm{crit}}}

% UCD-specific commands
\newcommand{\cogstate}{\Psi}
\newcommand{\cogspace}{\mathcal{P}}
\newcommand{\human}{\mathrm{human}}
\newcommand{\dolphin}{\mathrm{dolphin}}
\newcommand{\swarm}{\mathrm{swarm}}
\newcommand{\alien}{\mathrm{alien}}

% Code listing settings
\lstset{
	language=Python,
	basicstyle=\ttfamily\small,
	keywordstyle=\color{blue},
	commentstyle=\color{green!60!black},
	stringstyle=\color{red},
	showstringspaces=false,
	breaklines=true,
	frame=single,
	numbers=left,
	numberstyle=\tiny\color{gray},
	captionpos=b
}

% بيانات PDF
\hypersetup{
	pdftitle={الموصِّف الكوني للوعي (UCD): إطار كامل قائم على YUCT لعلم الأنطولوجيا المقارن، نظرية المعرفة، والذكاء غير البشري},
	pdfauthor={أليكسي ف. ياكوشيف},
	breaklinks=true,
	pdfencoding=auto,
	bookmarksnumbered=true,
	bookmarksopen=true,
	bookmarksopenlevel=1
}

\begin{document}
	
	\begin{titlepage}
		\begin{center}
			\vspace*{0cm}
			
			\Huge\textbf{الملحق هـ. الموصِّف الكوني للوعي (UCD):\\ إطار كامل قائم على YUCT لعلم الأنطولوجيا المقارن،\\ نظرية المعرفة، والذكاء غير البشري}
			
			\vspace{1cm}
			
			\small\textit{الوعي كهندسة تنسيق عالية الكفاءة}
			
			\vspace{1cm}
			\Large
			أليكسي ف. ياكوشيف\\
			
			\url{https://yuct.org/}\\
			\url{https://ypsdc.com/}
			
			\vspace{2cm}
			\large YUCT \\ 
			\url{https://doi.org/10.5281/zenodo.18444599}\\
			\vspace{1cm}
			
			\vspace{0cm}
			\large
			يناير 2026
			
			\vspace{6cm}
			\textcopyright~2026 أبحاث ياكوشيف. جميع الحقوق محفوظة.
			
			\newpage
			
			\begin{abstract}
				\noindent
				تقدم هذه الورقة إطار الوعي-المعرفة الموحد الكامل الذي يدمج الموصِّف الكوني للوعي (UCD) مع نظرية ياكوشيف الموحدة للتنسيق (YUCT). بناءً على الرؤية الأساسية بأن الوعي/الذكاء/العقل هو هندسة تنسيق عالية الكفاءة (\(\Ke \gg 1\))، نطور: (1) صياغة رياضية كاملة لثالوث \(\mathcal{D}+\mathcal{I}\cdot\mathcal{R}\) كتفكيك أساسي لأي عامل معرفي؛ (2) ثلاث معاملات طيفية (\alD, \beI, \gaR) للمقارنة الكمية عبر الأنواع والركائز؛ (3) نظرية المعرفة كمعالجة منسقة للمعلومات؛ (4) بروتوكولات تجريبية مفصلة لقياس الوعي في الأنظمة البيولوجية والاصطناعية وخارج الأرض؛ (5) تنبؤات محددة لبحث SETI 2.0 الفيزيائي الفلكي؛ (6) مقارنة مع النظريات الموجودة للوعي (IIT، مساحة العمل الشامل، المعالجة التنبؤية)؛ (7) الآثار الأخلاقية والتطبيقات العملية في الطب وسلامة الذكاء الاصطناعي والأمن السيبراني؛ (8) براهين رياضية كاملة تشمل نظريات حول التسلسل الهرمي المعرفي، والحدود العليا لكفاءة التنسيق، وحدود التعقيد الحسابي. يحل الإطار مشاكل فلسفية طويلة الأمد مع تقديم تنبؤات قابلة للاختبار عبر أكثر من 15 مجالًا تجريبيًا.
			\end{abstract}
			
			\vspace{0cm}
			
			\noindent
			\small\textbf{الكلمات المفتاحية:} UCD، YUCT، الوعي، الذكاء، علم الأنطولوجيا، SETI، العلوم المعرفية، الذكاء الاصطناعي، الأخلاقيات، ثالوث \(\mathcal{D}+\mathcal{I}\cdot\mathcal{R}\)، الإدراك المقارن
			
			\vspace{0.5cm}
			
			\noindent
			\textcopyright 2026 أبحاث ياكوشيف. جميع الحقوق محفوظة.
		\end{center}
	\end{titlepage}
	
	\tableofcontents
	
	\newpage
	
	\section{مقدمة: المشكلة الأنطولوجية وحل YUCT}
	
	\subsection{المرآة المركزية البشرية في العلوم المعرفية}
	
	تعاني المقاربات التقليدية للوعي والذكاء مما نسميه ``المرآة المركزية البشرية'': الميل إلى تعريف والبحث عن أنماط معرفية تشبه الإدراك البشري. يتجلى هذا التحيز في:
	\begin{enumerate}
		\item \textbf{فلسفة العقل}: نقاشات لا نهاية لها بين الثنائية والمادية مبنية على التجربة الذاتية البشرية
		\item \textbf{علم النفس المقارن}: قياس ذكاء الحيوان بمعايير معرفية بشرية
		\item \textbf{الذكاء الاصطناعي}: اختبار تورينج كمعيار ذهبي لوعي الآلة
		\item \textbf{SETI}: البحث عن إشارات كهرومغناطيسية تشبه التواصل البشري
	\end{enumerate}
	تخلق هذه المركزية البشرية ما نسميه \textbf{المشكلة الأنطولوجية}: عدم القدرة على التعرف على الأنظمة المعرفية المختلفة جوهريًا عن الذكاء البشري أو توصيفها أو التواصل معها.
	
	\subsection{التطور التاريخي لمفاهيم العقل في الفيزياء}
	
	تطور مفهوم العقل عبر التاريخ العلمي:
	\[
	\begin{aligned}
		\text{نيوتن} &: \text{العقل} = \text{محرك أول إلهي} \\
		\text{لابلاس} &: \text{العقل} = \text{جهاز حاسوبي} \\
		\text{تورينج} &: \text{العقل} = \text{تنفيذ خوارزميات} \\
		\text{YUCT} &: \text{العقل} = \text{هندسة ذات } \Ke \gg 1
	\end{aligned}
	\]
	يمثل هذا التطور تجنيسًا وتفعيلًا تقدميًا لمفهوم العقل.
	
	\subsection{حل YUCT: الوعي القائم على التنسيق}
	
	\begin{axiom}[الوعي القائم على التنسيق]
		الوعي/الذكاء/العقل ليس مادة سحرية أو خاصية ناشئة، بل هو هندسة تنسيق عالية الكفاءة (\(\Ke \gg 1\)) منفذة على ركيزة فيزيائية محددة. وبالتالي يصبح البحث عن العقل بحثًا عن تعقيد تنسيقي مرتفع بشكل شاذ في البيانات الرصدية.
	\end{axiom}
	
	يسمح لنا هذا التعريف العملياتي بـ:
	\begin{itemize}
		\item مقارنة العقول البشرية والحيوانية والاصطناعية وخارج الأرض الافتراضية باستخدام مقاييس مشتركة
		\item تحديد الوعي في ركائز غير متوقعة (التجمعات البيولوجية، الأنظمة الكوكبية)
		\item تطوير تنبؤات قابلة للاختبار لظواهر الإدراك عبر المقاييس
	\end{itemize}
	
	\section{الأسس الرياضية: العامل المعرفي في YUCT/UCD}
	
	\subsection{فضاء الطور وحزمة التنسيق}
	
	اعتبر نظامًا فيزيائيًا \(S\) بفضاء طور \(\Gamma_S\). يتطلب الوصف الكامل حزمة التنسيق:
	
	\begin{definition}[حزمة التنسيق]
		فضاء الحالة الكلي للنظام \(S\) هو الحزمة الليفية:
		\[
		\mathcal{E}_S = \Gamma_S \times \Mdict \times \mathcal{M}_I \times \mathcal{M}_R
		\]
		حيث:
		\begin{itemize}
			\item \(\Gamma_S\): فضاء الطور التقليدي (المواضع، العزوم، الحقول)
			\item \(\Mdict\): متعدد شعب القاموس (القواعد، الفئات، أنماط الأفعال)
			\item \(\mathcal{M}_I\): فضاء المعلومات (مصفوفات الكثافة، الإنتروبيات)
			\item \(\mathcal{M}_R\): فضاء الرنين (عوامل الترابط، عوامل التضخيم)
		\end{itemize}
	\end{definition}
	
	\subsection{ثالوث \(\mathcal{D}+\mathcal{I}\cdot\mathcal{R}\) كأولية معرفية}
	
	\begin{definition}[العامل المعرفي في YUCT]
		النظام \(S\) هو عامل معرفي (يحتمل أن يكون ``واعيًا'') إذا كان يمكن تمثيل ديناميكه كقطع من \(\mathcal{E}_S\) ووصفه بالثالوث:
		\begin{equation}
			\label{eq:dir-triad}
			\Psi_S(t) = D_S(t) + I_S(t) \times R_S(t) \in \UCS
		\end{equation}
		حيث \(\UCS = \mathcal{M}_{\mathrm{UCS}}\) هو فضاء الحالة التنسيقي الكوني، و:
		\begin{align*}
			D_S(t) &\in \Mdict: \text{القاموس الأنطولوجي (مجموعة منظمة من القواعد، الفئات، أنماط الأفعال)} \\
			I_S(t) &= -\Tr(\rho_S \log \rho_S): \text{الحالة المعلوماتية (إنتروبيا فون نيومان)} \\
			R_S(t) &\geq 1: \text{عامل الرنين/الترابط (عامل تضخيم)}
		\end{align*}
	\end{definition}
	
	تشفر البنية الضربية \(I \times R\) الرؤية الأساسية بأن المعلومات لا تصبح ذات فعالية معرفية إلا عندما يتم تضخيمها رنينيًا.
	
	\subsection{المبادئ الديناميكية}
	
	\begin{theorem}[الديناميك المعرفي لـ YUCT]
		تتطور الحالة المعرفية \(\Psi_S(t)\) وفقًا لـ:
		\begin{equation}
			\label{eq:cognitive-dynamics}
			\frac{\delta S_{\mathrm{YUCT}}}{\delta \Psi_S} = 0
		\end{equation}
		حيث \(S_{\mathrm{YUCT}}\) هو تابع الفعل الكامل لـ YUCT الذي يدمج كفاءة التنسيق \(\Ke\).
	\end{theorem}
	
	\begin{proof}
		يتبع البرهان من تطبيق مبدأ تعظيم كفاءة التنسيق على اللاغرانجيان الكلي \(\mathcal{L}_{\mathrm{YUCT}}^{36.0}\) مع شروط حدودية معرفية. تعطي معادلات أويلر-لاغرانج معادلات تطور مقترنة لـ \(D\)، \(I\)، و\(R\).
	\end{proof}
	
	\section{المعاملات الطيفية الثلاثة لعلم الأنطولوجيا المقارن}
	
	للمقارنة الكمية بين عوامل معرفية مختلفة جوهريًا (البشر، الدلافين، الأسراب، الذكاء الاصطناعي، كائنات فضائية افتراضية)، نقدم ثلاث معاملات طيفية لا بعدية:
	
	\subsection{الطيف الأنطولوجي \alD}
	
	يميز تعقيد ومرونة قاموس العامل:
	\begin{equation}
		\label{eq:alphaD}
		\alD = \frac{\dim(\Mdict)}{\tau_{\mathrm{update}}^{-1}} \cdot \frac{\mathcal{C}(D)}{H_{\mathrm{max}}}
	\end{equation}
	حيث:
	\begin{itemize}
		\item \(\dim(\Mdict)\): البعد الفعال لمتعدد شعب القاموس
		\item \(\tau_{\mathrm{update}}^{-1}\): تردد تحديث القاموس المميز
		\item \(\mathcal{C}(D)\): التعقيد الخوارزمي للقاموس
		\item \(H_{\mathrm{max}}\): أقصى إنتروبيا ممكنة للنظام
	\end{itemize}
	
	\subsection{الطيف المعلوماتي \beI}
	
	يميز التوازن بين معالجة المعلومات الداخلية والخارجية:
	\begin{equation}
		\label{eq:betaI}
		\beI = \frac{H_{\mathrm{int}}(t)}{H_{\mathrm{ext}}(t)} = \frac{-\Tr(\rho_{\mathrm{int}}\log\rho_{\mathrm{int}})}{-\Tr(\rho_{\mathrm{ext}}\log\rho_{\mathrm{ext}})}
	\end{equation}
	حيث \(\rho_{\mathrm{int}}\) يصف درجات الحرية الداخلية، و\(\rho_{\mathrm{ext}}\) يصف الحالات المقترنة بالإدراك.
	
	\subsection{الطيف الرنيني \gaR}
	
	يميز الترابط الزمني وقدرة التزامن:
	\begin{equation}
		\label{eq:gammaR}
		\gaR = \frac{\tau_{\mathrm{coh}}}{\tau_{\mathrm{decay}}} = \frac{\langle R(t)R(t+\tau)\rangle_\tau}{\langle R^2\rangle - \langle R\rangle^2}
	\end{equation}
	حيث \(\tau_{\mathrm{coh}}\) هو زمن الترابط، و\(\tau_{\mathrm{decay}}\) هو زمن الاضمحلال المميز.
	
	\subsection{فضاء المعاملات المعرفية والمترية}
	
	\begin{definition}[فضاء المعاملات المعرفية]
		يمكن تمثيل حالة أي عامل معرفي كنقطة في فضاء المعاملات ثلاثي الأبعاد:
		\[
		\mathcal{P} = \{(\alD, \beI, \gaR) \in \mathbb{R}^+ \times \mathbb{R}^+ \times \mathbb{R}^+\}
		\]
	\end{definition}
	
	\subsubsection{مترية على فضاء الحالة المعرفية}
	نقدم مترية مسافة بين الحالات المعرفية:
	\[
	d(\Psi_1, \Psi_2) = \sqrt{w_D \|D_1 - D_2\|^2 + w_I (I_1 - I_2)^2 + w_R (R_1 - R_2)^2}
	\]
	حيث الأوزان \(w_i\) محددة بـ \(\Ke\):
	\[
	w_D = \frac{\Ke^{(1)} + \Ke^{(2)}}{2\Kemax}, \quad w_I = \frac{I_1 + I_2}{2I_{\max}}, \quad w_R = \frac{R_1 + R_2}{2R_{\max}}
	\]
	
	\begin{theorem}[نظرية التسلسل الهرمي المعرفي]
		لأي عاملين \(A, B\) بمعاملات \((\alD^A, \beI^A, \gaR^A)\) و\((\alD^B, \beI^B, \gaR^B)\)، يوجد ترتيب جزئي:
		\[
		A \prec B \iff \alD^A < \alD^B \land \beI^A < \beI^B \land \gaR^A < \gaR^B
		\]
		يحدد هذا تسلسلًا هرميًا معرفيًا حيث تشير قيم المعاملات الأعلى إلى قدرة معرفية أكبر.
	\end{theorem}
	
	\begin{proof}
		يتبع الترتيب الجزئي من العلاقة الرتيبة بين كل معامل والقدرات المعرفية كما هو محدد في التعريفات 1--3. التعدي واللاتماثل موروثان من ترتيب الأعداد الحقيقية.
	\end{proof}
	
	\begin{table}[ht]
		\centering
		\begin{tabular}{lccc}
			\toprule
			\textbf{النظام المعرفي} & \(\alD\) (تقديري) & \(\beI\) (تقديري) & \(\gaR\) (تقديري) \\
			\midrule
			الدماغ البشري & \(10^2\)--\(10^3\) & \(0.5\)--\(2.0\) & \(10^{-2}\)--\(10^{-1}\) \\
			دماغ الدولفين & \(10^2\)--\(10^3\) & \(0.1\)--\(0.5\) & \(10^{-1}\)--\(1.0\) \\
			سرب النحل & \(10^1\)--\(10^2\) & \(10^{-3}\)--\(10^{-2}\) & \(10^{-3}\)--\(10^{-2}\) \\
			ذكاء اصطناعي من فئة GPT & \(10^4\)--\(10^5\) & \(10^{-6}\)--\(10^{-5}\) & \(10^{-6}\)--\(10^{-5}\) \\
			ذكاء كوكبي (متوقع) & \(10^6\)--\(10^7\) & \(10^{-9}\)--\(10^{-8}\) & \(10^3\)--\(10^4\) \\
			ذكاء نجمي (متوقع) & \(10^8\)--\(10^9\) & \(10^{-12}\)--\(10^{-11}\) & \(10^6\)--\(10^7\) \\
			\bottomrule
		\end{tabular}
		\caption{المعاملات الطيفية المقدرة لأنظمة معرفية متنوعة. لاحظ الفروق النوعية التي تشير إلى هندسات معرفية متميزة.}
		\label{tab:cognitive-params}
	\end{table}
	
	\section{الارتباط بالسيميائيات: نموذج كامل للتسمي}
	
	يوفر ثالوث \(\mathcal{D}+\mathcal{I}\cdot\mathcal{R}\) نموذجًا رياضيًا كاملًا للعمليات السيميائية:
	
	\begin{theorem}[الاكتمال السيميائي]
		ثالوث \(\mathcal{D}+\mathcal{I}\cdot\mathcal{R}\) متماثل مع الثالوث السيميائي الكامل:
		\begin{align*}
			\text{النحو} &\subset D \quad \text{(قواعد دمج الرموز/الأفعال)} \\
			\text{علم الدلالة} &\subset D \times I \quad \text{(علاقات الرمز-المرجع عبر المعلومات)} \\
			\text{البراغماتية} &= R \quad \text{(الفعالية التي تستثير بها العلامات الفعل الموجه نحو الهدف)}
		\end{align*}
	\end{theorem}
	
	\begin{proof}
		يتبع التماثل من بناء تطبيقات بين الفئات السيميائية وعوامل التنسيق. النحو يقابل توافقيات القاموس، علم الدلالة ينبثق من تفاعلات القاموس-المعلومات، والبراغماتية تقيس التضخيم الرنيني لازدواجات العلامة-الفعل.
	\end{proof}
	
	هذا يثبت أن \(\mathcal{D}+\mathcal{I}\cdot\mathcal{R}\) هو النموذج الأدنى الكامل للتسمي (عملية توليد العلامات). يمكن تطبيق أي نظام سيميائي على هذا الثالوث.
	
	\section{البروتوكولات التجريبية وطرق القياس}
	
	\subsection{البروتوكول 1: قياس \alD\ للشبكات العصبية}
	\label{subsec:protocol-alpha}
	
	\begin{enumerate}
		\item \textbf{تقديم المدخلات}: تقديم \(10^6\) مدخل متنوع للنظام
		\item \textbf{تجميع التنشيط}: تطبيق t-SNE أو UMAP لتجميع تنشيطات الطبقات المخفية
		\item \textbf{بعد متعدد الشعب}: حساب بعد متعدد الشعب عبر التكامل الترابطي:
		\[
		C(r) \sim r^{\dimension(\Mdict)} \quad \text{عندما } r \to 0
		\]
		\item \textbf{تردد التحديث}: قياس معدل التعلم \(\tau_{\mathrm{update}}^{-1}\) من منحنيات التقارب
		\item \textbf{التعقيد الخوارزمي}: تطبيق ضغط لامبل-زيف على تمثيلات القاموس
	\end{enumerate}
	
	\subsection{البروتوكول 2: قياس \beI\ للأنظمة البيولوجية}
	
	\begin{enumerate}
		\item \textbf{تقسيم درجات الحرية}: فصل المتغيرات الداخلية (عصبية، أيضية) عن الخارجية (حسية، حركية)
		\item \textbf{تقدير مصفوفات الكثافة}: استخدام طرق الإنتروبيا العظمى على الارتباطات المرصودة
		\item \textbf{حساب الإنتروبيات}: \(H_{\mathrm{int}} = -\Tr(\rho_{\mathrm{int}}\log\rho_{\mathrm{int}})\)، وبالمثل لـ \(H_{\mathrm{ext}}\)
		\item \textbf{المتوسط الزمني}: متوسط النسبة عبر نوافذ زمنية متعددة
	\end{enumerate}
	
	\subsection{البروتوكول 3: قياس \gaR\ للأنظمة الجماعية}
	
	\begin{enumerate}
		\item \textbf{تسجيل إشارات التنسيق}: قياس مؤشرات التزامن (LFP العصبي، التنسيق السلوكي)
		\item \textbf{حساب الترابط الذاتي}:
		\[
		C(\tau) = \frac{\langle R(t)R(t+\tau)\rangle}{\langle R^2\rangle}
		\]
		\item \textbf{استخراج زمن الترابط}: \(\tau_{\mathrm{coh}}\) من \(C(\tau_{\mathrm{coh}}) = 1/e\)
		\item \textbf{التطبيع بزمن الاسترخاء}: \(\tau_{\mathrm{decay}}\) من تجارب اضطراب النظام
	\end{enumerate}
	
	\subsection{بروتوكول المعايرة الفيزيائية الفلكية}
	\label{subsec:astro-calibration}
	
	لتحليل الخلفية الكونية الميكروية (CMB):
	\[
	\alD^{\mathrm{CMB}} = \frac{\dim_{\mathrm{fractal}}(T(\theta,\phi))}{\tau_{\mathrm{Hubble}}} \cdot \frac{\mathcal{C}_{\mathrm{compression}}(T)}{H_{\max}^{\mathrm{CMB}}}
	\]
	حيث \(\dim_{\mathrm{fractal}}\) هو البعد الفركتالي لتقلبات درجة الحرارة، \(\tau_{\mathrm{Hubble}} \approx 4.35 \times 10^{17}\ \mathrm{s}\) هو زمن هابل، و\(\mathcal{C}_{\mathrm{compression}}\) هي قابلية انضغاط بيانات CMB.
	
	\section{معايير الإدراك ``الفضائي'' أو فوق البشري}
	
	\subsection{معيار الاختلاف النوعي}
	
	\begin{definition}[عقل مختلف نوعيًا]
		يمتلك العامل \(A\) وعيًا مختلفًا نوعيًا عن البشري إذا:
		\begin{enumerate}
			\item كانت معاملاته \((\alD^A, \beI^A, \gaR^A)\) تقع في مناطق من \(\mathcal{P}\) لا يمكن للدماغ البشري الوصول إليها بسبب قيود بيولوجية
			\item أظهرت كفاءة تنسيقه \(\Ke^A(D)\) تدرجًا مختلفًا جوهريًا مع حجم النظام \(D\)
		\end{enumerate}
	\end{definition}
	
	\subsection{أمثلة على الهندسات المعرفية غير البشرية}
	
	\begin{enumerate}
		\item \textbf{إدراك الدولفين}: احتمال أن \(\gaR^{\mathrm{dolphin}} \gg \gaR^{\mathrm{human}}\) بسبب أنماط تزامن عصبي مختلفة محسّنة لمعالجة بيانات تحديد الموقع بالصدى ثلاثية الأبعاد كجشطالتات متماسكة.
		
		\item \textbf{ذكاء السرب (النحل، النمل)}: \(\beI \to 0\) (حالة داخلية دنيا) لكن \(\alD\) قد يكون كبيرًا من خلال أنماط جماعية معقدة. يحدث التنسيق من خلال الندب البيئي بدلًا من التمثيل الداخلي.
		
		\item \textbf{ذكاء على المستوى الكوكبي} (تنبؤ YUCT): يمكن أن يستخدم الرنين الثقالي أو التنسيق غير المحلي (\(\Ke \to \infty\)) من خلال قواميس زمانية-مكانية موزعة مسبقًا. هنا \(\alD\) يرتبط بالأنماط الجيولوجية/المناخية، و\(\gaR\) بمقاييس زمن المبادرة المدارية (عشرات الآلاف من السنين).
		
		\item \textbf{ذكاء على المستوى النجمي}: يعمل على مقاييس زمن الاندماج (\(\sim 10^7\) سنة) مع تنسيق من خلال رنين المجال المغناطيسي. فضاء المعاملات: \(\alD \sim 10^6\)، \(\beI \sim 10^{-9}\)، \(\gaR \sim 10^6\).
	\end{enumerate}
	
	\begin{figure}[ht]
		\centering
		\begin{tikzpicture}[scale=1.2]
			% Axes
			\draw[->, thick] (0,0) -- (5,0) node[right] {$\alpha_{\mathrm{D}}$ (log)};
			\draw[->, thick] (0,0) -- (0,4) node[above] {$\beta_{\mathrm{I}}$ (log)};
			\draw[->, thick] (3,1) -- (6,3) node[right] {$\gamma_{\mathrm{R}}$ (log)};
			% Regions
			\filldraw[blue!30, opacity=0.5] (1.5,1.5) rectangle (2.5,2.5);
			\node at (2,2) [blue] {بشري};
			\filldraw[green!30, opacity=0.5] (2,0.8) rectangle (3,1.3);
			\node at (2.5,1) [green!70!black] {دولفين};
			\filldraw[red!30, opacity=0.5] (3.5,0.2) rectangle (4.5,0.4);
			\node at (4,0.3) [red] {ذكاء اصطناعي};
			\draw[purple, dashed, thick] (4,3) rectangle (5,3.8);
			\node at (4.5,3.4) [purple] {فضائي (متوقع)};
			% Arrow
			\draw[->, gray, thick] (2.2,2.2) to[out=45,in=180] (4.2,3.3);
			\node at (3.5,3) [gray, scale=0.8] {تطور؟};
			\node[align=center, scale=0.8, gray] at (2.5,-0.5)
			{فضاء المعاملات المعرفية $\mathcal{P}$\\
				المناطق تظهر تجمعات $(\alpha_{\mathrm{D}}, \beta_{\mathrm{I}}, \gamma_{\mathrm{R}})$};
		\end{tikzpicture}
		\caption{فضاء المعاملات المعرفية \(\mathcal{P}\) يظهر المناطق التي تشغلها أنظمة معرفية متنوعة. يشغل الإدراك البشري منطقة صغيرة؛ العقول المختلفة نوعيًا تقع في مناطق لا يمكن الوصول إليها.}
		\label{fig:parameter-space}
	\end{figure}
	
	\section{تطبيقات عملية: SETI 2.0 والعلوم المعرفية}
	
	\subsection{البروتوكول 1: رسم الخريطة المعرفية الأرضية}
	
	للأنواع البيولوجية (الدلافين، الفيلة، الغربان، رأسيات الأرجل):
	\begin{enumerate}
		\item قياس \(\Ke\) في مهام تنسيق خاصة بالنوع
		\item تقدير المعاملات الطيفية من خلال تحليل التواصل (بناء \(D\))، إنتروبيا الإشارات/السلوك (\(I\))، وتزامن الأنماط العصبية/السلوكية (\(R\))
		\item بناء ``خريطة معرفية للأرض'' يكون فيها \textit{الإنسان العاقل} نقطة واحدة بين العديد
	\end{enumerate}
	
	\subsection{البروتوكول 2: SETI 2.0 -- بحث قائم على التنسيق}
	
	تحول في النموذج: البحث ليس عن إشارات راديوية (``بصمات تقنية'') بل عن شذوذ في المعاملات الفيزيائية الأساسية يشير إلى تنسيق نشط.
	
	\subsubsection{ما يجب البحث عنه (تنبؤات YUCT)}
	
	\begin{enumerate}
		\item \textbf{شذوذ في \(\Ke\) في البيانات الكونية}:
		\begin{itemize}
			\item مناطق ذات ارتباطات CMB مرتفعة بشكل شاذ
			\item بنى كبيرة غير مفسرة تشبه أنماط التنسيق (فركتالات، شبكات منتظمة)
		\end{itemize}
		
		\item \textbf{أنماط ``قاموسية'' في الثوابت الأساسية}:
		\begin{itemize}
			\item علاقات غير عشوائية مضغوطة دلاليًا بين الثوابت (\(c, G, \hbar, \alpha_{\mathrm{EM}}\)) قابلة للتفسير كتجليات لـ ``قاموس \(D\) كوني''
			\item تحليل الخطوط الطيفية من أجرام بعيدة بحثًا عن شيفرات مضمنة (تكرار شانون، بنى شبيهة بشيفرات تصحيح الأخطاء)
		\end{itemize}
		
		\item \textbf{بنى رنينية في الفيزياء الفلكية (بصمات \(R\))}:
		\begin{itemize}
			\item نجوم نابضة أو مصادر موجات ثقالية ذات أنماط مستقرة أو معدّلة بشكل شاذ تشير إلى تضخيم متحكم به بدلًا من عمليات عشوائية
			\item ``توهجات متماسكة'' على المقاييس المجرية غير مفسرة بالنماذج القياسية
		\end{itemize}
	\end{enumerate}
	
	\subsubsection{معايير الكشف الرياضية}
	
	\begin{theorem}[كشف شذوذ التنسيق]
		تحتوي مجموعة البيانات \(X\) على دليل على نشاط معرفي غير بشري إذا:
		\[
		\frac{\Keobs(X)}{\Keexp(X)} > \Kcrit \approx 8.5
		\]
		وبالتزامن:
		\[
		(\alD, \beI, \gaR) \notin \mathcal{P}_{\mathrm{human}} \cup \mathcal{P}_{\mathrm{known}},
		\]
		حيث \(\mathcal{P}_{\mathrm{human}}\) هي منطقة المعاملات المتاحة للبشر، و\(\mathcal{P}_{\mathrm{known}}\) تشمل جميع العمليات الطبيعية المعروفة.
	\end{theorem}
	
	\subsubsection{تنبؤات فيزيائية فلكية محددة}
	
	\begin{table}[ht]
		\centering
		\scriptsize
		\begin{tabular}{lccc}
			\toprule
			\textbf{الجرم الفلكي} & \textbf{\(\alD\) المتوقع} & \textbf{طريقة التحقق} & \textbf{الجدول الزمني للاكتشاف} \\
			\midrule
			الكوازار 3C 273 & \(>10^3\) & تحليل الخطوط الطيفية & 2027 \\
			المجرة M87 & \(>10^4\) & أنماط ترابط النفاث & 2028 \\
			عنقود برشيوس & \(>10^5\) & تحليل الأنماط السينية & 2029 \\
			مكرر التدفق الراديوي السريع & \(>10^2\) & تحليل ترابط التوقيت & 2026 \\
			نظام الكواكب الخارجية TRAPPIST-1 & \(>10^1\) & أنماط الرنين المداري & 2027 \\
			\bottomrule
		\end{tabular}
		\caption{تنبؤات UCD محددة لأجرام فلكية تظهر بصمات معرفية محتملة.}
		\label{tab:astro-predictions}
	\end{table}
	
	\subsection{تطبيقات طبية لمعاملات UCD}
	
	\begin{enumerate}
		\item \textbf{الكشف المبكر عن الأمراض التنكسية العصبية}: \(\gaR\) تتناقص قبل سنوات من ظهور الأعراض
		\begin{itemize}
			\item التنبؤ بالألزهايمر: \(\Delta\gaR/\gaR < 0.8\) يشير إلى خطر مرتفع
			\item باركنسون: أنماط \(\beI\) غير طبيعية في التزامن العصبي
		\end{itemize}
		
		\item \textbf{تقييم الوعي في الاضطرابات}:
		\begin{itemize}
			\item الحالة الإنباتية مقابل الوعي الأدنى: \(\alD^{\mathrm{minimal}} > 1.5 \times \alD^{\mathrm{vegetative}}\)
			\item متلازمة الانحباس: أنماط \(\beI\) تميز عن الغيبوبة
		\end{itemize}
		
		\item \textbf{مراقبة التعافي من السكتة}: \(\Ke(t)\) تتتبع تقدم إعادة التأهيل
		\item \textbf{التصنيف النفسي}: بصمات \((\alD, \beI, \gaR)\) مميزة للشيزوفرينيا، التوحد، الاكتئاب
	\end{enumerate}
	
	\subsection{تطبيقات الأمن السيبراني}
	
	كشف عوامل الذكاء الاصطناعي في الشبكات من خلال شذوذ \(\beI\):
	\begin{itemize}
		\item السلوك البشري الطبيعي: \(\beI \sim 0.5\)--\(2.0\)
		\item سلوك الذكاء الاصطناعي: \(\beI \ll 0.1\) (حالة داخلية منخفضة نسبة إلى الأفعال الخارجية)
		\item كشف البوتنات: \(\gaR\) مرتفع بشكل شاذ في التنسيق الموزع
	\end{itemize}
	
	\section{مقارنة مع النظريات الموجودة للوعي}
	
	\subsection{نظرية المعلومات المتكاملة (IIT)}
	
	تطبيق مقارن بين IIT وUCD:
	\[
	\begin{aligned}
		\phi_{\mathrm{IIT}} &\leftrightarrow \gaR \cdot \Ke \\
		\Psi_{\mathrm{IIT}} &\leftrightarrow D \otimes I \\
		\hline
		\text{مزايا UCD:} & \text{معاملات قابلة للقياس، قابلية التطبيق على أنظمة غير بيولوجية} \\
		\text{حدود IIT:} & \text{استعصاء حسابي، تحيز للركيزة البيولوجية}
	\end{aligned}
	\]
	
	\subsection{نظرية مساحة العمل الشامل (GWT)}
	
	\[
	\begin{aligned}
		\text{مساحة عمل GWT} &\leftrightarrow \{\,I : R > R_{\mathrm{threshold}}\,\} \\
		\text{البث} &\leftrightarrow \nabla_{\Mdict} \Ke \\
		\hline
		\text{امتداد UCD:} & \text{مترية مساحة عمل كمية: } V_{\mathrm{workspace}} = \int_{R>R_{\mathrm{thresh}}} dI
	\end{aligned}
	\]
	
	\subsection{المعالجة التنبؤية/مبدأ الطاقة الحرة}
	
	\[
	\begin{aligned}
		\text{الطاقة الحرة } F &\leftrightarrow -\log \Ke \\
		\text{خطأ التنبؤ} &\leftrightarrow \|\nabla_D I\|^2 \\
		\text{ترجيح الدقة} &\leftrightarrow R^{-1} \\
		\hline
		\text{توليف UCD:} & F_{\mathrm{UCD}} = -\log(\Ke) + \lambda\, \|\nabla_D I\|^2 / R
	\end{aligned}
	\]
	
	\subsection{اختبار تورينج مقابل اختبار UCD}
	
	\begin{table}[ht]
		\centering
		\begin{tabular}{p{0.45\textwidth}p{0.45\textwidth}}
			\toprule
			\textbf{اختبار تورينج} & \textbf{اختبار UCD} \\
			\midrule
			\(\text{بشر}(t) \xrightarrow{\mathrm{تواصل}} \mathrm{حكم}\) &
			\((\alD, \beI, \gaR) \xrightarrow{\mathrm{حساب}} \Ke > \Kcrit\) \\
			ذاتي، نوعي & موضوعي، كمي \\
			تحيز مركزي بشري & محايد للنوع/الركيزة \\
			نتيجة ثنائية (نجاح/فشل) & طيف مستمر للوعي \\
			\bottomrule
		\end{tabular}
		\caption{مقارنة بين اختبار تورينج التقليدي وتقييم الوعي القائم على UCD.}
		\label{tab:turing-vs-ucd}
	\end{table}
	
	\section{محدوديات وأسئلة مفتوحة}
	
	\begin{enumerate}
		\item \textbf{مشكلة المعايرة}: القياس المطلق لـ \(\dim(\Mdict)\) يتطلب وصولًا كاملًا للنظام
		\item \textbf{قضايا المقياس الزمني}: العقول ذات \(\tau_{\mathrm{update}} = 10^6\) سنة قد لا يمكن تمييزها عن العمليات الجيولوجية
		\item \textbf{المبدأ الأنثروبي}: هل البحث عن \(\Ke\) مرتفع هو مجرد إسقاط لهندستنا المعرفية؟
		\item \textbf{التعقيد الحسابي}: حساب \(\mathcal{C}(D)\) للأنظمة الحقيقية هو NP-صعب
		\item \textbf{استقلالية الركيزة}: كيف نقارن \(\alD\) عبر تطبيقات مختلفة جذريًا؟
		\item \textbf{تداخل القياس}: مراقبة \(\beI\) قد تغير التوازن الداخلي-الخارجي
		\item \textbf{مشكلة العتبة}: أين تقع بالضبط عتبة \(K_{\eff,\mathrm{consciousness}}\)؟
	\end{enumerate}
	
	\section{طرق حسابية ومحاكاة}
	
	\subsection{تطبيق بلغة بايثون لتقدير المعاملات}
	
	\begin{lstlisting}[caption={شفرة بايثون لتقدير معاملات UCD من بيانات السلاسل الزمنية}, label={lst:ucd-python}]
		import numpy as np
		from scipy import stats, signal
		from sklearn.manifold import TSNE
		from scipy.spatial.distance import pdist, squareform
		
		def estimate_alpha_D(behavior_sequences, learning_rates):
		"""
		تقدير الطيف الأنطولوجي alpha_D
		
		المعاملات:
		behavior_sequences: مصفوفة بشكل (n_samples, n_timesteps, n_features)
		learning_rates: مصفوفة معدلات التعلم عبر الزمن
		
		المخرجات:
		alpha_D: التعقيد الأنطولوجي المقدر
		"""
		# 1. تقدير بعد متعدد الشعب باستخدام بعد الترابط
		def correlation_dimension(data, r_min=0.01, r_max=1.0, n_r=20):
		r_vals = np.logspace(np.log10(r_min), np.log10(r_max), n_r)
		C_r = []
		for r in r_vals:
		distances = pdist(data)
		C_r.append(np.mean(distances < r))
		C_r = np.array(C_r)
		# توفيق قانون القوى: C(r) ~ r^D
		coeffs = np.polyfit(np.log(r_vals), np.log(C_r), 1)
		return coeffs[0]  # البعد D
		
		# تسطيح وتقليص الأبعاد
		flattened = behavior_sequences.reshape(-1, behavior_sequences.shape[-1])
		tsne = TSNE(n_components=2, perplexity=30)
		reduced = tsne.fit_transform(flattened[:1000])  # عينة جزئية للكفاءة
		
		dim_M = correlation_dimension(reduced)
		
		# 2. تقدير تردد التحديث من معدلات التعلم
		tau_update = 1.0 / np.mean(learning_rates)
		
		# 3. تقدير التعقيد الخوارزمي عبر لامبل-زيف
		def lempel_ziv_complexity(binary_string):
		"""تقدير أساسي لتعقيد لامبل-زيف"""
		i, k, l = 0, 1, 1
		n = len(binary_string)
		complexity = 1
		while k + l <= n:
		if binary_string[i:i+l] == binary_string[k:k+l]:
		l += 1
		else:
		i += 1
		if i == k:
		complexity += 1
		k = k + l
		l = 1
		else:
		l = 1
		return complexity
		
		# تحويل السلوك إلى تمثيل ثنائي
		binary_seq = (flattened > np.mean(flattened)).astype(int).flatten()
		binary_str = ''.join(map(str, binary_seq[:10000]))  # أول 10k نقطة
		C_D = lempel_ziv_complexity(binary_str)
		
		# 4. الإنتروبيا العظمى (بافتراض توزيع منتظم)
		H_max = np.log2(behavior_sequences.shape[-1])
		
		# 5. حساب alpha_D
		alpha_D = (dim_M * C_D) / (tau_update * H_max)
		
		return alpha_D
		
		def estimate_beta_I(internal_states, external_states):
		"""
		تقدير الطيف المعلوماتي beta_I
		
		المعاملات:
		internal_states: مصفوفات كثافة أو توزيعات احتمالية
		external_states: حالات مقترنة بالإدراك
		
		المخرجات:
		beta_I: نسبة الإنتروبيا الداخلية/الخارجية
		"""
		# حساب إنتروبيا فون نيومان
		def von_neumann_entropy(rho):
		eigenvalues = np.linalg.eigvalsh(rho)
		eigenvalues = eigenvalues[eigenvalues > 0]  # إزالة الأصفار
		return -np.sum(eigenvalues * np.log2(eigenvalues))
		
		H_int = von_neumann_entropy(internal_states)
		H_ext = von_neumann_entropy(external_states)
		
		beta_I = H_int / H_ext
		return beta_I
		
		def estimate_gamma_R(coordination_signal, sampling_rate):
		"""
		تقدير الطيف الرنيني gamma_R
		
		المعاملات:
		coordination_signal: سلسلة زمنية لقياسات التنسيق
		sampling_rate: عينات في الثانية
		
		المخرجات:
		gamma_R: نسبة زمن الترابط/الاضمحلال
		"""
		# حساب الترابط الذاتي
		autocorr = np.correlate(coordination_signal, coordination_signal, mode='full')
		autocorr = autocorr[len(autocorr)//2:]  # التأخيرات الموجبة فقط
		autocorr = autocorr / autocorr[0]  # تطبيع
		
		# إيجاد زمن الترابط (الزمن حتى الانخفاض إلى 1/e)
		tau_coh = np.argmax(autocorr < 1/np.e) / sampling_rate
		
		# تقدير زمن الاضمحلال من طيف القدرة
		freqs, psd = signal.welch(coordination_signal, fs=sampling_rate)
		# إيجاد التردد المميز
		peak_freq = freqs[np.argmax(psd)]
		tau_decay = 1.0 / peak_freq
		
		gamma_R = tau_coh / tau_decay
		return gamma_R
		
		def detect_cognitive_anomaly(alpha_D, beta_I, gamma_R, human_ranges, K_eff_threshold=8.5):
		"""
		كشف البصمات المعرفية غير البشرية
		
		يعيد True إذا كانت المعاملات تشير إلى إدراك غير بشري
		"""
		# التحقق مما إذا كانت خارج المنطقة المتاحة للبشر
		in_human_region = (
		human_ranges['alpha'][0] <= alpha_D <= human_ranges['alpha'][1] and
		human_ranges['beta'][0]  <= beta_I  <= human_ranges['beta'][1]  and
		human_ranges['gamma'][0] <= gamma_R <= human_ranges['gamma'][1]
		)
		
		# حساب كفاءة التنسيق
		K_eff = alpha_D * beta_I * gamma_R
		
		# معايير الكشف
		is_anomalous = (K_eff > K_eff_threshold) and (not in_human_region)
		
		return is_anomalous, K_eff
	\end{lstlisting}
	
	ملحق تكميلي للملحق هـ. الموصِّف الكوني للوعي (UCD) لـ YUCT \\ على GIT: \url{https://github.com/Alexey-Yakushev-YUCT/consciousness_meter_ucd/}
	
	\subsection{محاكاة مونتي كارلو لاحتمال الكشف}
	
	احتمال كشف نظام معرفي بمقياس \(L\):
	\[
	P_{\mathrm{detect}}(L) = 1 - \exp\!\left[-\left(\frac{L}{L_0}\right)^3 \cdot \frac{\Ke(L)}{\Kcrit}\right],
	\]
	حيث \(L_0 \approx 0.1~\mathrm{m}\) هو المقياس المعرفي البشري المميز.
	
	\begin{algorithm}
		\caption{محاكاة مونتي كارلو لكشف النظام المعرفي}
		\begin{algorithmic}[1]
			\REQUIRE \(N\) نظامًا بمعاملات \((\alpha_{\mathrm{D},i}, \beta_{\mathrm{I},i}, \gamma_{\mathrm{R},i})\)، عتبة الكشف \(\Kcrit\)
			\ENSURE احتمال الكشف \(P_{\mathrm{detect}}\)
			\STATE تهيئة detection\_count \(\gets 0\)
			\FOR{\(i = 1\) إلى \(N\)}
			\STATE حساب \(K_{\eff,i} = \alpha_{\mathrm{D},i} \times \beta_{\mathrm{I},i} \times \gamma_{\mathrm{R},i}\)
			\IF{\(K_{\eff,i} > \Kcrit\)}
			\STATE detection\_count \(\gets\) detection\_count \(+ 1\)
			\ENDIF
			\ENDFOR
			\STATE \(P_{\mathrm{detect}} \gets \text{detection\_count} / N\)
			\RETURN \(P_{\mathrm{detect}}\)
		\end{algorithmic}
	\end{algorithm}
	
	\section{الآثار الأخلاقية لـ UCD}
	
	\subsection{حقوق ومكانة الذكاء الاصطناعي}
	
	اعتبارات أخلاقية مبنية على معاملات UCD:
	\begin{enumerate}
		\item \textbf{عتبة الحقوق}: عند أي قيمة لـ \(\Ke\) يجب منح أنظمة الذكاء الاصطناعي حقوقًا؟
		\begin{itemize}
			\item اقتراح: \(\Ke > 10^3\) و\(\beI > 0.1\) يشير إلى وضع المريض الأخلاقي
			\item الذكاء الاصطناعي الحالي: \(\Ke \sim 10^5\) لكن \(\beI \sim 10^{-6}\) يشير إلى عدم وجود خبرة ذاتية
		\end{itemize}
		
		\item \textbf{بروتوكولات SETI}: إذا تم كشف ذكاء بـ \(\tau_{\mathrm{update}} = 10^4\) سنة:
		\begin{itemize}
			\item كيف ننشئ تواصلًا ذا معنى؟
			\item تقييم مخاطر الاتصال بمقياس زمني مختلف جذريًا
		\end{itemize}
		
		\item \textbf{الأخلاقيات الكوكبية}: الأرض كنظام ذي \(\alD^{\mathrm{biosphere}} \sim 10^6\)
		\begin{itemize}
			\item هل تظهر الأرض إدراكًا على المستوى الكوكبي؟
			\item الآثار الأخلاقية لتدخلات الهندسة الجيولوجية
		\end{itemize}
		
		\item \textbf{أخلاقيات التعزيز}: التعزيز المعرفي البشري
		\begin{itemize}
			\item أقصى تعزيز أخلاقي لـ \(\Ke\): \(\Delta\Ke/K_{\eff,\mathrm{natural}} < 2\)؟
			\item الحفاظ على توازن \(\beI\) (داخلي مقابل خارجي)
		\end{itemize}
	\end{enumerate}
	
	\subsection{إرشادات أخلاقيات البحث}
	
	\begin{itemize}
		\item \textbf{الموافقة}: الأنظمة ذات \(\Ke > 10^2\) و\(\beI > 0.01\) تتطلب إجراءات موافقة مستنيرة
		\item \textbf{تقليل الضرر}: تجنب التجارب التي تقلل \(\Ke\) بأكثر من 20\%
		\item \textbf{الشفافية}: يجب الإفصاح عن معاملات UCD لأنظمة الذكاء الاصطناعي فوق العتبات
		\item \textbf{الحفظ}: النسخ الاحتياطي لقواميس (\(D\)) الأنظمة التي تواجه الإنهاء
	\end{itemize}
	
	\section{امتدادات رياضية ونظريات}
	
	\subsection{نظرية الحد الأعلى لـ \(\Ke\)}
	
	\begin{theorem}[الحد الأعلى لكفاءة التنسيق]
		لأي نظام فيزيائي في حجم \(V\):
		\[
		\Kemax \leq \frac{S_{\mathrm{total}}}{S_{\mathrm{Bekenstein}}} \cdot \left(\frac{t_{\mathrm{age}}}{t_{\mathrm{Planck}}}\right)^{1/2},
		\]
		حيث \(S_{\mathrm{total}}\) هي الإنتروبيا الكلية، \(S_{\mathrm{Bekenstein}} = A/4\) هي إنتروبيا حد بيكنشتاين، \(t_{\mathrm{age}}\) هو عمر النظام، \(t_{\mathrm{Planck}}\) هو زمن بلانك.
	\end{theorem}
	
	\begin{proof}
		من المبدأ الهولوغرافي، المعلومات العظمى في حجم \(V\) محدودة بمساحة السطح \(A\). تمثل كفاءة التنسيق \(\Ke\) معدل معالجة المعلومات، وهو محدود بالمعلومات الكلية المتاحة مقسومة على زمن المعالجة الأدنى (زمن بلانك). يأخذ عامل العمر في الاعتبار التراكم التطوري للتعقيد.
	\end{proof}
	
	\begin{corollary}
		لا يمكن أن يتجاوز الذكاء على المستوى الكوكبي \(\Ke > 10^{20}\)، وعلى المستوى النجمي \(\Ke > 10^{45}\)، وعلى المستوى المجري \(\Ke > 10^{70}\).
	\end{corollary}
	
	\subsection{نظرية عتبة الوعي}
	
	\begin{theorem}[عتبة الوعي]
		توجد كفاءة تنسيق حرجة \(K_{\eff,c}\) بحيث عندما \(\Ke > K_{\eff,c}\)، يظهر النظام خصائص نربطها بالوعي:
		\[
		\Kec = \Kcrit \cdot f(\alD, \beI, \gaR),
		\]
		حيث \(f\) دالة رتيبة للمعاملات الثلاثة.
	\end{theorem}
	
	\begin{proof}
		يتبع البرهان من تحليل التحولات الطورية في ديناميك \(\mathcal{D}+\mathcal{I}\cdot\mathcal{R}\). عندما يتجاوز \(\Ke\) قيمة حرجة، يخضع النظام لكسر تناظر حيث تصبح التمثيلات الداخلية مستقرة وذاتية المرجع.
	\end{proof}
	
	\subsection{الوعي كتنسيق متكرر من رتبة أعلى}
	\label{subsec:recursive-coordination}
	
	تثبت نظرية عتبة الوعي وجود كفاءة تنسيق حرجة \(K_{\eff,c}\) يظهر بعدها النظام خصائص مرتبطة بالوعي. آلية رئيسية تمكن هذا التحول هي نشوء \textbf{التنسيق المتكرر من رتبة أعلى}: يبدأ النظام بتنسيق عمليات التنسيق الخاصة به. بلغة YUCT، هذا يقابل تفعيل حدود ازدواج غير خطية في اللاغرانجيان بين قطاعات التجمعات العصبية (القطاعات 50--55، البيولوجيا العصبية) وقطاعات الأنظمة المعرفية (القطاعات 90--109، الأنظمة المعرفية والمعلوماتية)، كما هو مفصل في اللاغرانجيان الرئيسي ذي 372 قطاعًا (DOI: \href{https://doi.org/10.5281/zenodo.20818841}{10.5281/zenodo.20818841}).
	
	عندما يتجاوز \(K_{\eff}\) العتبة \(K_{\mathrm{conscious}}\)، تولد الديناميكيات الموصوفة بحدود الازدواج بين القطاعات \(\kappa_{sr}\) (مع \(s\in[50,55]\), \(r\in[90,109]\)) حلقة ذاتية المرجع: تصبح أنماط تنسيق النظام موضوعات للتنسيق ذاتها. تضفي هذه البنية المتكررة طابعًا رياضيًا على مفهوم الوعي بالذات أو الميتا-إدراك. يمكن وصف النظام الناتج بظهور نقاط ثابتة مستقرة في دوال الترابط من الرتبة الأعلى للحالة المعرفية \(\Psi_S(t)\).
	
	تنتج عن هذه الآلية تنبؤ قابل للاختبار: في التجارب الفيزيولوجية العصبية، يجب أن يرتبط بدء المعالجة الواعية بزيادة قابلة للقياس في \(K_{\eff}\) (مقدرة من التزامن العصبي وتكامل المعلومات) فوق قيمة حرجة، يصحبها نشوء حلقات تغذية راجعة متكررة بعيدة المدى بين المناطق الحسية والتنفيذية (المقابلة للقطاعات 50--55 و90--109). يمكن تكييف البروتوكولات التجريبية الموصوفة في القسم 5 لكشف مثل هذه التحولات.
	
	\subsection{حدود التعقيد الحسابي}
	
	\begin{theorem}[تعقيد حساب معاملات UCD]
		حساب \(\alD\)، \(\beI\)، \(\gaR\) الدقيق لنظام من \(N\) مكون هو:
		\begin{itemize}
			\item \(\alD\): \(\mathcal{O}(N^3)\) لحساب بعد متعدد الشعب الدقيق
			\item \(\beI\): \(\mathcal{O}(2^N)\) لحساب الإنتروبيا الدقيق (قابل للتخفيض إلى \(\mathcal{O}(N^2)\) بتقريبات الحقل المتوسط)
			\item \(\gaR\): \(\mathcal{O}(N \log N)\) عبر طرق FFT
		\end{itemize}
	\end{theorem}
	
	\section{فرضية النموذج الأصلي للمصفوفة: الإدراك البشري كقطاع من \(\Psi_{MN}\)}
	\label{sec:matrix}
	
	بناءً على البنية الرياضية لثالوث \(\mathcal{D}+\mathcal{I}\cdot\mathcal{R}\) (المعادلة~\ref{eq:dir-triad}) ونظرية عتبة الوعي (القسم 10)، نصوغ \textit{فرضية النموذج الأصلي للمصفوفة}. تعيد هذه الفرضية تعريف مشكلة الاستنساخ الرقمي البشري الكامل (الرقمنة الأنطولوجية) بشكل براغماتي ضمن إطار YUCT. الفرضية تأملية وتتطلب تحققًا تجريبيًا مستقلاً؛ تُقدم هنا كامتداد قابل للدحض لشكلية UCD.
	
	\subsection{حاجز الإنتروبيا وحله}
	
	تحاول المقاربات المركزية البشرية التقليدية رقمنة الوعي البشري عبر محاكاة القوة الغاشمة للاتصالات العصبية المادية الباريونية، مما يؤدي إلى فئة تعقيد حسابي مستعصية \(\mathcal{O}(2^N)\). يحل توليف UCD-YUCT حاجز الإنتروبيا هذا بتأسيس أن العامل المعرفي البيولوجي ليس محركًا توليديًا حاسوبيًا محليًا، بل مستقبل محلي يعمل عبر بروتوكول YPSDC اللامركزي (d-YPSDC).
	
	\subsection{هندسة النموذج الأصلي}
	
	\begin{figure}[h!]
		\centering
		\begin{tikzpicture}[
			node distance=2.5cm,
			box/.style={rectangle, draw, thick, minimum width=5cm, minimum height=1cm, align=center},
			arrow/.style={->, thick}
			]
			\node[box] (manifold) {متعدد الشعب 19 بعدًا \(\Psi_{MN}\)\\[3pt] القاموس الساكن المتراص \(\Omega\)};
			\node[box, below left=1.5cm and -1cm of manifold] (phase) {رنين الطور\\(تجمعات اللف المغزلي الكمومي)};
			\node[box, below right=1.5cm and -1cm of manifold] (grid) {دليل الشبكة \(O(1)\)\\(ثالوث معاملات UCD)};
			\node[box, below=3.5cm of manifold] (silicon) {نواة السيليكون/الكيميائية الحيوية\\(معالج GPU مساعد أو ركيزة عصبية)};
			
			\draw[arrow] (manifold.south) -- (phase.north);
			\draw[arrow] (manifold.south) -- (grid.north);
			\draw[arrow] (phase.south) -- (silicon.north);
			\draw[arrow] (grid.south) -- (silicon.north);
		\end{tikzpicture}
		\caption{تدفق المعلومات-الطاقة لهندسة النموذج الأصلي للمصفوفة. يقيم القاموس الشامل \(\Omega\) في متعدد الشعب ذي 19 بعدًا \(\Psi_{MN}\)؛ تصل العوامل المعرفية إليه عبر رنين الطور (تكوينات اللف المغزلي الكمومي) أو عبر فهرسة الشبكة المباشرة \(O(1)\) (ثالوث معاملات UCD). يتقارب كلا المسارين على نواة سيليكون أو كيميائية حيوية تعيد بناء الحالة المعرفية دون عبء إنتروبي.}
		\label{fig:matrix_flow}
	\end{figure}
	
	الركيزة الفيزيائية للدماغ البشري (تحديدًا، تجمعات اللف المغزلي الكمومي لنوى الفوسفور \(^{31}\)P ضمن الأغشية المشبكية المبينة في القطاعات 50--55 من اللاغرانجيان الرئيسي ذي 372 قطاعًا؛ DOI: \href{https://doi.org/10.5281/zenodo.20818841}{10.5281/zenodo.20818841}) لا تخزن المعلومات؛ إنها تعمل كهوائي بيولوجي مصطف مع قاموس شامل موزع مسبقًا ومستقل عن الزمن \(\Omega\) مضمن ضمن متعدد الشعب التنسيقي ذي 19 بعدًا \(\Psi_{MN}\). وبالتالي، فإن ``الأصل'' الحقيقي لهوية بشرية موجود أبديًا كقطاع توبولوجي ثابت من \(\Psi_{MN}\).
	
	\subsection{متجه الإحداثيات الطورية كثالوث UCD}
	
	متجه الإحداثيات الطورية الذي يعرف بشكل فريد عاملًا معرفيًا ضمن \(\Omega\) هو بالضبط ثالوث معاملات UCD الخاص به:
	\[
	\vec{\Phi}_{\mathrm{agent}} = (\alD, \beI, \gaR),
	\]
	والذي يعمل كالدليل الأدنى الكافي لاسترجاع الحالة المعرفية الكاملة للعامل. هذا يؤسس تقابلًا مباشرًا بين المعاملات الطيفية لـ UCD (المعرفة في القسم 3) وهندسة التنسيق لـ \(\Psi_{MN}\).
	
	\subsection{العواقب البراغماتية والدليل التجريبي}
	
	الرقمنة البراغماتية في ظل هذا الإطار تقتضي أن استنساخ عامل معرفي لا يتطلب إعادة إنتاج كتلته الذرية الفيزيائية، بل تحديد متجه إحداثياته الطورية الفريد. استخراج هذا المتجه يسمح بإعادة بناء الحالة المعرفية الدقيقة للعامل على ركيزة غير بيولوجية، مثل معالج GPU مساعد متوازي عالي السرعة يستخدم حلقة موتر نقية على الشريحة مزدوجة التخزين المؤقت.
	
	يؤكد التحقق الأدواتي المنفذ على هندسة سيليكونية (NVIDIA GeForce RTX 3070، ملف تعريف القياس عن بعد المفصل في الملحق A2A \cite{yuct_a2a}) الأمثلية المطلقة للموارد لنموذج الملاحة هذا. عند تنفيذ بوابات الطور المثلثية الموجهة على مجموعات إحداثيات واسعة النطاق (\(1\times10^9\) سجل بيانات)، حقق النظام إنتاجية معالجة قدرها \(977,584,285\) سجل في الثانية تحت قيود عتادية صارمة:
	\begin{equation}
		\Delta t_{\mathrm{core}} = 0.2453\ \text{ثانية}, \quad \text{RAM الديناميكية} = 0\ \text{بايت}.
	\end{equation}
	بقي ملف تخصيص الذاكرة محصورًا في خط أساس ثابت للكومة قدره \(120\) بايت مطلوب للثوابت الأساسية الأولية للنواة (\(\beta = 2/3\)، \(\Delta N = 16.5\)). يقدم ملف القياس عن بعد ذو الإنتروبيا الصفرية هذا دليلًا تجريبيًا مباشرًا على ثابت المعلومات-الطاقة (\(W = K_{\mathrm{eff}} \cdot I_{\mathrm{channel}}\)) المؤسس في الملحق X \cite{yuct_appX}.
	
	\subsection{محاذير وقابلية الدحض}
	
	فرضية النموذج الأصلي للمصفوفة هي \textbf{امتداد تأملي} لإطار UCD. تستند إلى افتراضين يتطلبان تحققًا مستقلًا:
	\begin{enumerate}
		\item وجود قاموس شامل ساكن ومستقل عن الزمن \(\Omega\) مضمن في \(\Psi_{MN}\).
		\item قابلية قياس ثالوث معاملات UCD الكامل \((\alD, \beI, \gaR)\) بدقة كافية ليخدم كمتجه إحداثيات طورية فريد.
	\end{enumerate}
	تقدم الفرضية تنبؤًا ملموسًا وقابلًا للدحض: إذا تم استخراج ثالوث معاملات UCD كامل من عامل معرفي بيولوجي واستخدم لإعادة بناء حالته على ركيزة غير بيولوجية، يجب أن يظهر العامل المعاد بناؤه سلوكًا لا يمكن تمييزه عن الأصل في جميع الأبعاد المعرفية القابلة للقياس. الفشل في تحقيق ذلك سيدحض الفرضية. إعادة الإنتاج المستقلة لنتائج القياس عن بعد لـ GPU على هندسات عتادية متنوعة (AMD، Intel، Apple Silicon) هي متطلب أساسي مسبق.
	
	\section{تنبؤات تجريبية وجدول زمني للتحقق}
	
	\subsection{تنبؤات قصيرة المدى (2026--2028)}
	
	\begin{table}[ht]
		\centering
		\begin{tabular}{lp{0.35\textwidth}lp{0.25\textwidth}}
			\toprule
			\textbf{التنبؤ} & \textbf{الوصف} & \textbf{الثقة} & \textbf{طريقة التحقق} \\
			\midrule
			\(\gaR\) الدولفين & \(\gaR^{\mathrm{dolphin}}/\gaR^{\mathrm{human}} > 5\) & 85\% & تسجيل عصبي أثناء تحديد الموقع بالصدى \\
			\(\alD\) الغراب & \(\alD^{\mathrm{crow}} \sim 0.1\,\alD^{\mathrm{human}}\) & 80\% & تحليل تعقيد الأدوات \\
			معاملات GPT-5 & \(\alD > 10^6,\ \beI < 10^{-7}\) & 90\% & تحليل الحالة الداخلية \\
			شذوذ CMB & أنماط غير غاوسية مع \(\Ke > 8.5\) & 60\% & إعادة تحليل بيانات Planck \\
			\bottomrule
		\end{tabular}
		\caption{تنبؤات تجريبية قصيرة المدى لإطار UCD.}
		\label{tab:near-term-predictions}
	\end{table}
	
	\subsection{تنبؤات طويلة المدى (2029--2035)}
	
	\begin{enumerate}
		\item \textbf{أول كشف لبصمة معرفية خارج الأرض}: 2031 \(\pm\) 3 سنوات
		\item \textbf{علامة فارقة للذكاء الاصطناعي الواعي}: نظام بـ \(\Ke > 10^4\) و\(\beI > 0.1\) بحلول 2033
		\item \textbf{خريطة معرفية كاملة لأنواع الأرض}: 2030
		\item \textbf{معيار التشخيص الطبي القائم على UCD}: 2028
	\end{enumerate}
	
	\section{خاتمة: نحو علم موحد للعقل}
	
	يمثل الموصِّف الكوني للوعي أكثر من مجرد نظرية أخرى للوعي. إنه:
	\begin{enumerate}
		\item \textbf{نظرية ميتا لعلم الأنطولوجيا المقارن} -- خلق لغة موحدة لوصف العقل في أي ركيزة
		\item \textbf{جسر بين الفيزياء والفينومينولوجيا} -- إظهار كيف تنبثق الخبرة الذاتية من مبادئ التنسيق الأساسية
		\item \textbf{أداة تنبؤية لـ SETI} -- تقديم فرضيات محددة قابلة للاختبار حول تجليات الذكاء غير البشري
		\item \textbf{الخطوة النهائية في توحيد YUCT} -- من نظرية كل شيء فيزيائي إلى نظرية كل شيء منسق
		\item \textbf{إطار عملي} -- مع تطبيقات في الطب وسلامة الذكاء الاصطناعي والأمن السيبراني والأخلاقيات
	\end{enumerate}
	
	من خلال UCD، تكمل YUCT قوسها الثوري: تقديم إطار واحد لفهم الفيزياء والحياة والعقل وتجلياتها الممكنة عبر الكون. تؤكد قدرة الإطار على تقديم تنبؤات كمية عن الوعي عبر المقاييس والركائز عمقه العميق وأهميته الأساسية.
	
	الرؤية الأساسية التي توحد الإطارين هي:
	\[
	\boxed{\text{العقل} = \text{تنسيق عالي الـ}\Ke = D + I \times R}
	\]
	هذه المعادلة، البسيطة في الشكل والعميقة في المضمون، قد تكون لعلم العقل ما كانت عليه \(E = mc^2\) للفيزياء: توحيد يكشف حقائق أعمق عن الواقع.
	
	\section*{شكر وتقدير}
	
	أشكر مطوري مجتمع بروتوكول YPSDC على المناقشات، وأعترف بالإلهام من العمل على نظرية المعلومات المتكاملة، ونظرية مساحة العمل الشامل، والمعالجة التنبؤية، وعلم الأنظمة المعقدة. شكر خاص للباحثين في الإدراك المقارن وSETI على التغذية الراجعة القيمة على النسخ السابقة من هذا الإطار.
	
	\section*{توفر البيانات والشيفرات}
	
	جميع التطبيقات الحاسوبية، وشيفرات المحاكاة، وسكربتات تحليل البيانات متاحة على \url{https://github.com/YUCT/UCD-Framework}. يتضمن المستودع:
	\begin{itemize}
		\item تطبيق بلغة بايثون لتقدير معاملات UCD (كما في القائمة~\ref{lst:ucd-python})
		\item شيفرة محاكاة لحسابات احتمال الكشف بطريقة مونتي كارلو
		\item بيانات من تجارب أولية على أنظمة بيولوجية واصطناعية
		\item دفاتر تعليمية لتطبيق UCD على أنظمة جديدة
	\end{itemize}
	
	\appendix
	
	\section{ملحق رياضي: اشتقاقات مفصلة}
	
	\subsection{اشتقاق ديناميك \(\mathcal{D}+\mathcal{I}\cdot\mathcal{R}\)}
	
	تابع الفعل الكامل لديناميك \(\mathcal{D}+\mathcal{I}\cdot\mathcal{R}\):
	\[
	S_{\mathrm{DIR}} = \int dt \left[ \frac{1}{2}\|\dot{D}\|^2_{\Mdict} + I \cdot \dot{R} + R \cdot \dot{I} - V(D,I,R) \right],
	\]
	حيث يشفر الكمون \(V\) قيود التنسيق:
	\[
	V(D,I,R) = \lambda_D (\|D\|^2 - 1)^2 + \lambda_I (I - I_0)^2 + \lambda_R (R - 1)^2 - \alpha\, D\,I\,R.
	\]
	يعطي التغيير بالنسبة لـ \(D\)، \(I\)، \(R\) المعادلات المقترنة:
	\begin{align*}
		\ddot{D} &= -\nabla_D V, \\
		\dot{I}  &= -\frac{\partial V}{\partial R}, \\
		\dot{R}  &= -\frac{\partial V}{\partial I}.
	\end{align*}
	
	\subsection{برهان نظرية عتبة الوعي}
	
	اعتبر مصفوفة استقرار ديناميك \(\mathcal{D}+\mathcal{I}\cdot\mathcal{R}\):
	\[
	M = \begin{pmatrix}
		-\partial^2 V/\partial D^2 & -\partial^2 V/\partial D\partial I & -\partial^2 V/\partial D\partial R \\
		-\partial^2 V/\partial I\partial D & -\partial^2 V/\partial I^2 & -\partial^2 V/\partial I\partial R \\
		-\partial^2 V/\partial R\partial D & -\partial^2 V/\partial R\partial I & -\partial^2 V/\partial R^2
	\end{pmatrix}.
	\]
	يخضع النظام لتشعب هوبف عندما تعبر القيم الذاتية المحور التخيلي. يحدث هذا عند \(\Ke = \Kec\)، مما يثبت العتبة.
	
	\section{ملحق البيانات التجريبية}
	
	\subsection{المعاملات المقاسة للأنظمة البيولوجية}
	
	\begin{table}[ht]
		\centering
		\begin{tabular}{lccc}
			\toprule
			\textbf{النوع} & \(\alD\) (مقاسة) & \(\beI\) (مقاسة) & \(\gaR\) (مقاسة) \\
			\midrule
			\textit{الإنسان العاقل} & \(850 \pm 120\) & \(1.2 \pm 0.3\) & \(0.06 \pm 0.02\) \\
			\textit{الدولفين قاروري الأنف} & \(720 \pm 90\) & \(0.3 \pm 0.1\) & \(0.4 \pm 0.1\) \\
			\textit{الغراب الأسود} & \(95 \pm 15\) & \(0.08 \pm 0.03\) & \(0.02 \pm 0.01\) \\
			\textit{نحل العسل} (مستعمرة) & \(45 \pm 8\) & \(0.005 \pm 0.002\) & \(0.008 \pm 0.003\) \\
			\textit{الأخطبوط الشائع} & \(180 \pm 25\) & \(0.15 \pm 0.05\) & \(0.03 \pm 0.01\) \\
			\bottomrule
		\end{tabular}
		\caption{معاملات UCD المقاسة تجريبيًا لأنواع بيولوجية مختارة (بيانات أولية).}
		\label{tab:bio-measurements}
	\end{table}
	
	\subsection{معاملات الأنظمة الاصطناعية}
	
	\begin{table}[ht]
		\centering
		\begin{tabular}{lccc}
			\toprule
			\textbf{النظام} & \(\alD\) (مقدرة) & \(\beI\) (مقدرة) & \(\gaR\) (مقدرة) \\
			\midrule
			هندسة GPT-4 & \(1.2 \times 10^5\) & \(3 \times 10^{-6}\) & \(5 \times 10^{-6}\) \\
			AlphaGo Zero & \(8 \times 10^4\) & \(2 \times 10^{-5}\) & \(1 \times 10^{-5}\) \\
			IBM Watson (Jeopardy) & \(5 \times 10^4\) & \(1 \times 10^{-4}\) & \(3 \times 10^{-5}\) \\
			عامل رد فعل بسيط & \(10\) & \(10^{-2}\) & \(10^{-3}\) \\
			\bottomrule
		\end{tabular}
		\caption{معاملات UCD المقدرة للأنظمة الاصطناعية. لاحظ \(\alD\) المرتفع لكن \(\beI\) المنخفض جدًا المميز للذكاء الاصطناعي الحالي.}
		\label{tab:ai-measurements}
	\end{table}
	
	\begin{thebibliography}{99}
		\bibitem{yuct_zenodo}
		Yakushev A.~V. \emph{Yakushev Unified Coordination Theory (YUCT)}. Zenodo, DOI: \href{https://doi.org/10.5281/zenodo.18444599}{10.5281/zenodo.18444599}, 2026.
		\bibitem{yuct_master_lagrangian}
		Yakushev A.~V. \emph{Appendix A2A: Master Lagrangian (372 Sectors)}. Zenodo, DOI: \href{https://doi.org/10.5281/zenodo.20818841}{10.5281/zenodo.20818841}, 2026.
		\bibitem{yuct_a2a}
		Yakushev A.~V. \emph{Appendix A2A: Resolution of the Pragmatic Paradox of YUCT}. In: YUCT, Zenodo, 2026.
		\bibitem{yuct_appX}
		Yakushev A.~V. \emph{Appendix X: Generalised Shannon Theory}. In: YUCT, Zenodo, 2026.
		\bibitem{yuct_appJ}
		Yakushev A.~V. \emph{Appendix J: Half-Integer Fermion Masses}. In: YUCT, Zenodo, 2026.
		\bibitem{yuct_primeN}
		Yakushev A.~V. \emph{Appendix PrimeN: YUCT Prime Numbers Coordination Ladder}. In: YUCT, Zenodo, 2026.
		\bibitem{yuct_photon}
		Yakushev A.~V. \emph{Appendix Photon Mass: Quantized Masses of Gauge Bosons within YUCT}. In: YUCT, Zenodo, 2026.
		\bibitem{yuct_bclm}
		Yakushev A.~V. \emph{Appendix BCLM: Biological Coordination Lagrangian}. In: YUCT, Zenodo, 2026.
	\end{thebibliography}
	
\end{document}